\documentclass[a4paper]{scrartcl}

\usepackage{csquotes}
\usepackage{amsmath}
\usepackage{amssymb}
\usepackage{amsthm}
\usepackage{graphicx}
\usepackage{float}
\usepackage{amsfonts}
\usepackage{subcaption}
\usepackage{url}
\usepackage{hyperref}
\usepackage{cleveref}
\usepackage{graphicx}
\usepackage[labelfont=bf]{caption}
\usepackage{bbm}

\bibliographystyle{vancouver}

% #1: path to the image
% #2: caption for the figure
% #3: label for the figure
% #4: optional height (default 7cm)
\newcommand{\insertfigure}[4][7cm]{%
    \begin{figure}[H]
        \centering
        \includegraphics[height=#1, width=\linewidth, keepaspectratio]{#2}
        \caption{#3}
        \label{#4}
    \end{figure}
}


% #1: path to the first image
% #2: path to the second image
% #3: caption for the entire figure
% #4: label for the figure
\newcommand{\insertfigures}[5][5cm]{
	\begin{figure}[H]
		\centering
		\begin{subfigure}[t]{0.45\textwidth}
			\centering
			\caption{}
			\includegraphics[height=#1, width=\linewidth,keepaspectratio]{#2}
		\end{subfigure}
		\hfill
		\begin{subfigure}[t]{0.45\textwidth}
			\centering
			\caption{}
			\includegraphics[height=#1, width=\linewidth, keepaspectratio]{#3}
		\end{subfigure}
		\caption{#4}
		\label{#5}
	\end{figure}
}

\pagenumbering{gobble}

\begin{document}

\begin{center}
    {
        Written report for \enquote{Learning from Distributed Data}\\
    }

    \vspace{1.5cm}

    {\large
        Report by\\
        Yannik Schroeder
    }

    \vspace{1.5cm}

    \today

    \vspace*{\fill}
\end{center}

\clearpage
\pagenumbering{Roman}

\tableofcontents

\pagebreak

\pagenumbering{arabic}

\section{Introduction}

Creating the training data

\begin{enumerate}
    \item df has columns \texttt{p1, p2, ..., is\_lv}
    \item X has columns \texttt{p1, p2, ..., p6}
    \item y has columns \texttt{is\_lv}
    \item \texttt{X\_train} and X\_test are a 80/20 random subset split of X
    \item y\_train, y\_test are a 80/20 random subset split of y. With \texttt{stratify=y} the ratio of true to false classes (which was about 3,000/100,000 = 3\%) is kept in y\_train and y\_test
    \item StandardScaler transforms X\_train to 0 mean and standard variance (1)
    \item It is fitted to the values from the training dataset and uses these to transform the test dataset. This avoids the test dataset to be too similar to the train dataset.
    \item X and y are split into N random subsets (cohesive regions would mostly contain 0's since there are so few 1's in the dataset)
    \item each client split contains the same number of 0 and 1 results, indices are random shuffled to avoid any regularity coming from original data
\end{enumerate}

\pagebreak

\addcontentsline{toc}{section}{References}

\bibliography{refs}

\end{document}